% !TEX root = ../../../main.tex
% Week 4, IMG_9236--IMG_9240. Corrections to tail indices and to the
% choice of an eventual geometric bound are recorded beside the proofs.
\section{Comparison tests}
\label{sec:analysis-i:comparison-tests}

\begin{definition}
  A series $\sum a_n$ is \emph{absolutely convergent} if
  $\sum |a_n|$ converges.
\end{definition}

\begin{theorem}[Comparison with a majorant]
  \label{thm:analysis-i:series-comparison}
  Suppose $|a_n|\le M_n$ for all sufficiently large $n$, where
  $M_n\ge0$ and $\sum M_n$ converges. Then $\sum |a_n|$ and
  $\sum a_n$ both converge.
\end{theorem}

\begin{proof}
  Let $T_n=\sum_{k=1}^n M_k$. For $m>n$ beyond the index where
  the comparison holds, the triangle inequality gives
  \[
    \left|\sum_{k=n+1}^{m}a_k\right|
    \le\sum_{k=n+1}^{m}|a_k|
    \le\sum_{k=n+1}^{m}M_k=T_m-T_n.
  \]
  Since $(T_n)$ converges, its Cauchy property makes the right-hand
  side smaller than any given $\varepsilon>0$ when $m>n$ are
  sufficiently large. The Cauchy criterion proves convergence of
  $\sum a_n$. Applying the same argument to $|a_n|$ proves absolute
  convergence as well.
\end{proof}
% IMG_9236 has T_m-T_(n+1). The sum from n+1 through m is T_m-T_n.

In particular, absolute convergence implies convergence by taking
$M_n=|a_n|$. A finite number of early terms has no effect on convergence,
because it only adds a fixed constant to all later partial sums.

\subsection{The \texorpdfstring{$p$}{p}-series}

For a real exponent $s>1$, consider $\sum_{n=1}^{\infty}1/n^s$.
The lecture uses real powers with their usual exponent and order laws;
their construction extends the rational powers developed earlier and
is deferred here.
% Clarification awaiting author review, concerning irrational exponents.
\begin{pendingreview}
The construction in \autoref{sec:analysis-i:roots-and-powers} defines
rational powers. For irrational $s$, the arguments here assume that
$n^s$ exists and obeys those order and exponent laws; they do not
construct it.
\end{pendingreview}
Group its terms into blocks between consecutive powers of $2$.
The block $2^j<n\le2^{j+1}$ has $2^j$ terms, and
\[
  \sum_{n=2^j+1}^{2^{j+1}}\frac{1}{n^s}
  \le 2^j\frac{1}{(2^j)^s}=2^{-j(s-1)}.
\]
Writing $q=2^{1-s}<1$, every partial sum is bounded by
\[
  1+\frac{1}{2^s}+\sum_{j=1}^{\infty}q^j<\infty.
\]
The positive partial sums increase, so the monotone bounded sequence
theorem proves convergence. When $s\le1$, we have $n^{-s}\ge1/n$
for $n\ge1$. Comparison with the harmonic partial sums proves
divergence. Thus the $p$-series converges exactly when $s>1$.

The lecture also states Euler's evaluation
$\sum_{n=1}^{\infty}1/n^2=\pi^2/6$, without proving its value.
The argument above proves convergence, not that evaluation.
\marginnote{The lecture asks how this relates to Gabriel's horn.
  Keep this question in mind for integration, where volume and surface
  area can be compared.}

\subsection{Exercises on comparison}

These exercises use estimates of positive terms and partial sums.
For the second, recall the eventual bounds supplied by the lower and
upper limits in \autoref{sec:analysis-i:upper-and-lower-limits}.

% !TEX root = ../../hw04.tex
% Source: HW4 Intro to Math Analysis I.pdf, page 1, question 4.
% The shift a is any positive real number, not just a positive integer.
\begin{exercise}
  \label{ex:hw04:shifted-root-series}
  \solutionlink{hw04:shifted-root-series}
  Prove that the telescope series
  \[
    \sum_{n=1}^{\infty}\left(\frac{1}{\sqrt n}-\frac{1}{\sqrt{n+a}}\right)
  \]
  converges for any $a>0$ and obtain the sum when $a=1$.
\end{exercise}

% !TEX root = ../../hw04.tex
% Source: HW4 Intro to Math Analysis I.pdf, page 1, question 5, parts (i)--(ii).
\begingroup
\renewcommand{\labelenumi}{(\roman{enumi})}
\begin{exercise}
  \label{ex:hw04:lower-upper-comparison}
  \solutionlink{hw04:lower-upper-comparison}
  Consider
  \[
    \sum_{n=1}^{\infty}a_n,\qquad\sum_{n=1}^{\infty}b_n,
  \]
  where $a_n\ge0$ and $b_n>0$ for $n=1,2,\ldots$. Define
  \[
    r=\liminf_{n\to\infty}\frac{a_n}{b_n},\qquad
    R=\limsup_{n\to\infty}\frac{a_n}{b_n}.
  \]
  \begin{enumerate}
    \item Show that if $r>0$ and $\sum_{n=1}^{\infty}b_n$ diverges,
      so does $\sum_{n=1}^{\infty}a_n$.
    \item Show that if $R<\infty$ and $\sum_{n=1}^{\infty}b_n$ converges,
      so does $\sum_{n=1}^{\infty}a_n$.
  \end{enumerate}
\end{exercise}
\endgroup


\subsection{The ratio test}

\begin{theorem}[Ratio test]
  \label{thm:analysis-i:ratio-test}
  Suppose $a_n\ne0$ for all sufficiently large $n$, and set
  $R_n=|a_{n+1}/a_n|$ on that tail.
  If $\limsup R_n<1$, then $\sum a_n$ converges absolutely.
  If $\liminf R_n>1$, then $\sum a_n$ diverges.
\end{theorem}

\begin{proof}
  First suppose $R_0=\limsup R_n<1$. Choose $\rho$ with
  $R_0<\rho<1$. By the eventual-bound property in
  \autoref{sec:analysis-i:upper-and-lower-limits}, some $N$ satisfies
  $R_n\le\rho$ for every $n\ge N$. Thus
  \[
    |a_{n+1}|\le\rho|a_n|\quad(n\ge N),\qquad
    |a_{N+k}|\le|a_N|\rho^k\quad(k\ge0),
  \]
  where the second inequality follows by induction. The geometric
  majorant $\sum_{k=0}^{\infty}|a_N|\rho^k$ converges, so
  \autoref{thm:analysis-i:series-comparison} proves absolute convergence.

  Now suppose $\liminf R_n>1$. Choose $\rho>1$ below that lower
  limit. For some $N$, all $n\ge N$ satisfy $R_n\ge\rho$, hence
  $|a_{N+k}|\ge|a_N|\rho^k$. Since $a_N\ne0$, these terms do not
  tend to zero. The series diverges by
  \autoref{thm:analysis-i:series-term-test}.
\end{proof}
% IMG_9239 uses R_0-epsilon as an upper bound. It must be above limsup,
% so the proof chooses R_0<rho<1. The divergence estimate is editorial
% explanation of the stated divergence implication.

\autoref{fig:geometric-bound-choice} highlights the choice in the first
part of the proof. The upper bound must exceed the limsup, while still
being small enough to produce a convergent geometric series.

\begin{marginfigure}
  \centering
  % !TEX root = ../../main.tex
% Shared physical dimensions for interval and number-line diagrams.
\tikzset{
  interval diagram/.style={
    x=1cm,y=6mm,font=\footnotesize,
    color=black,line width=0.5pt,line cap=round,line join=round,
    >={Stealth[length=4pt,width=3pt]}
  },
  interval endpoint/.style={
    circle,draw=black,line width=0.5pt,
    minimum size=4pt,inner sep=0pt,outer sep=0pt
  },
  interval open/.style={interval endpoint,fill=white},
  interval closed/.style={interval endpoint,fill=black},
  interval label/.style={inner sep=1pt}
}

\begin{tikzpicture}[interval diagram,x=38mm,y=6mm]
  \draw[gray,->] (0,0) -- (1.14,0);
  \draw (0.26,0) -- (1,0);
  \node[interval open] at (0.26,0) {};
  \node[interval closed] at (0.65,0) {};
  \node[interval open] at (1,0) {};
  \node[below=6pt] at (0.26,0) {$R_0$};
  \node[below=6pt] at (0.65,0) {$\rho$};
  \node[below=6pt] at (1,0) {$1$};
  \node[above=10pt] at (0.65,0) {$R_0<\rho<1$};
\end{tikzpicture}

  \caption[Choosing an eventual geometric bound.]{Choose $\rho$ in the open interval between $R_0=\limsup R_n$
    and $1$. The tail ratios are eventually below $\rho$, and
    $\sum\rho^k$ converges.}
  \label{fig:geometric-bound-choice}
\end{marginfigure}

For a geometric series with $x\ne0$, the ratios equal $|x|$.
The ratio test recovers absolute convergence when $|x|<1$.
The case $x=0$ is immediate without forming ratios.

\subsection{The root test}

\begin{theorem}[Root test]
  \label{thm:analysis-i:root-test}
  If $\limsup_{n\to\infty}|a_n|^{1/n}<1$, then $\sum a_n$
  converges absolutely.
\end{theorem}

\begin{proof}
  Choose $\rho$ strictly between the upper limit and $1$.
  All sufficiently late terms satisfy $|a_n|^{1/n}\le\rho$,
  so $|a_n|\le\rho^n$. The geometric majorant $\sum\rho^n$
  converges, and comparison proves the claim.
\end{proof}
% IMG_9240 gives only the convergence criterion and again writes
% R_0-epsilon. The same correction as in the ratio test applies.

A limiting value of $1$ does not decide convergence. For example,
both $1/n$ and $1/n^2$ have successive-term ratios tending to $1$
and $n$th-root limits equal to $1$, but only the second series converges.

The next exercise focuses on the divergence implications, using a lower
limit to bound the late terms away from zero.

% !TEX root = ../../hw04.tex
% Source: HW4 Intro to Math Analysis I.pdf, page 2, question 6, parts (i)--(ii).
% Source introduction: "In class, we discussed ratio and root tests for
% convergence. We now consider the tests for divergence."
% Replace the classroom framing with a mathematical introduction.
\begingroup
\renewcommand{\labelenumi}{(\roman{enumi})}
\begin{exercise}
  \label{ex:hw04:divergence-tests}
  \solutionlink{hw04:divergence-tests}
  Consider series
  \[
    \sum_{n=1}^{\infty}a_n,\qquad a_n>0,\quad n=1,2,\ldots.
  \]
  The following are ratio and root tests for divergence.
  \begin{enumerate}
    \item Define
      \[
        R_1=\liminf_{n\to\infty}\frac{a_{n+1}}{a_n}.
      \]
      Show that if $R_1>1$, then the series diverges.
    \item Define
      \[
        R_2=\liminf_{n\to\infty}(a_n)^{1/n}.
      \]
      Show that if $R_2>1$, then the series diverges.
  \end{enumerate}
\end{exercise}
\endgroup

